(1) 外转折点满足 \(E=A/b\)，故 \(b=A/E\)。在 \(R<r<b\)，经典径向动量为虚数，WKB 指数作用量定义为 \(S=(1/\hbar)\int_R^b\sqrt{2m(A/r-E)}\,dr\)，它是无量纲正数；衰变概率的指数是 \(-2S\) 而不是 \(-S\)。
(2) 设 \(\epsilon=R/b\in(0,1)\)，作 \(r=b\sin^2\theta\)，则根号内变成 \(\sqrt{2mE}\cot\theta\)，而 \(dr=2b\sin\theta\cos\theta\,d\theta\)。积分得到 \(S=(b\sqrt{2mE}/\hbar)[\arccos\sqrt\epsilon-\sqrt{\epsilon(1-\epsilon)}]\)，在 \(\epsilon\to1\) 时自然趋零。
(3) 在禁阻区宽、指数作用量 \(S\gg1\) 的 WKB 近似下，透射率的主指数为 \(\mathcal T\asymp e^{-2S}\)，预指数需由核内匹配与边界条件决定。若粒子对势垒每秒尝试频率为 \(\nu\)，粗略衰变常数 \(\lambda_d\simeq\nu\mathcal T\)，半衰期 \(t_{1/2}\simeq(\ln2)/(\nu\mathcal T)\)；这不是从库仑势单独确定的绝对衰变率。
(4) 当 \(\epsilon\ll1\)，\(\arccos\sqrt\epsilon-\sqrt{\epsilon(1-\epsilon)}=\pi/2-2\sqrt\epsilon+O(\epsilon^{3/2})\)。主导项 \(2S\simeq\pi A\sqrt{2m}/(\hbar\sqrt E)-4\sqrt{2mAR}/\hbar+\cdots\)；能量升高使外转折点靠近原子核，禁阻区缩短而穿透显著增强。若靠近转折点太近或势垒几乎消失，简单的大 \(S\) 指数近似失去精度。
